% To be reformatted in the journal style upon acceptance.
\documentclass[11pt]{amsart}

\usepackage{amsmath,amssymb}
\usepackage[hidelinks]{hyperref}

\newtheorem{thm}{Theorem}[section]
\newtheorem{lem}[thm]{Lemma}
\newtheorem{cor}[thm]{Corollary}
\newtheorem{prop}[thm]{Proposition}
\theoremstyle{definition}
\newtheorem{defi}[thm]{Definition}
\newtheorem{rem}[thm]{Remark}
\newtheorem{ques}[thm]{Question}
\newtheorem*{thmstar}{Theorem}

\newcommand{\DM}{\mathrm{DM}}
\newcommand{\cube}{\square_{\mathrm{DM}}}
\newcommand{\cubefp}{\square_{\mathrm{DM}}^{\mathrm{fp}}}
\newcommand{\Kar}{\mathrm{Kar}}
\newcommand{\lc}{Q}
\newcommand{\bnd}{\partial}
\newcommand{\Face}{\mathrm{Face}}
\newcommand{\Diag}{\mathrm{Diag}}
\newcommand{\rmin}{\mathrm{r}}
\newcommand{\Up}{\mathord{\uparrow}}

\begin{document}

\title[The strict layer of De Morgan cubical sets]{The strict layer
  of De Morgan cubical sets:\\ minimal fillers, a coherence
  dichotomy,\\ and the Karoubi envelope}

\author{Ryota Shibusawa}
\address{Daiichi Institute of Technology, Japan}
\email{r-shibusawa@daiichi-koudai.ac.jp}

\subjclass[2020]{06D30; 18N40, 18B25, 03B38, 55U35}
\keywords{cubical sets, De Morgan algebra, Eilenberg--Zilber
  category, Karoubi envelope, projective algebras, degeneracy,
  thin filler, cubical type theory}

\begin{abstract}
  The De Morgan cube category $\cube$ --- cubes with faces,
  degeneracies, symmetries, both connections, reversal, and
  diagonals, the shape category of CCHM cubical type theory --- is
  not an Eilenberg--Zilber category: we show it is not even
  idempotent complete, and identify its Karoubi envelope as the
  opposite of the category of finite \emph{projective} De Morgan
  algebras, which a theorem of Bova and Cabrer characterizes
  combinatorially --- a bridge between the cubical and the
  unification-theoretic literatures that appears not to have been
  noticed.  We then analyze the structure that survives the failure
  of Eilenberg--Zilber normal forms.  Building on a companion study
  of boundary fibers in free De Morgan algebras (fibers are Boolean
  intervals $2^{D}$ over antichains of diagonal points), we prove a
  one-step \emph{defect formula}: deleting the minimal diagonal
  points of any element reaches the least element of its boundary
  fiber, so the failure of any substitution to preserve minimality
  is measured by an explicit antichain.  Our main result is a
  \emph{coherence dichotomy}: for a substitution $f$ and an element
  $g$, correcting before or after transport agree if and only if the
  transported boundary avoids the defect antichain of $g$ --- an
  equational criterion that holds unconditionally on the wide
  subcategory generated by faces, symmetries and reversal, and fails
  in a controlled way for degeneracies, connections and diagonals;
  in particular the minimal skeleton is not even cylinder-stable.
  All statements are verified exhaustively by machine through
  dimension three ($7{,}828{,}354$ elements).  We conclude that in
  the absence of Eilenberg--Zilber structure, thinness of fillers
  cannot be a property and must be carried as structure --- the
  median interpolation of the companion paper being the minimal such
  structure.
\end{abstract}

\maketitle

\section{Introduction}
\label{sec:intro}

Simplicial sets owe much of their combinatorial control to the
Eilenberg--Zilber lemma: every cell is uniquely a degeneracy of a
nondegenerate cell.  Cube categories with connections but without
diagonals enjoy the same property \cite{Campion}; the shape category
of Cohen--Coquand--Huber--M\"ortberg cubical type theory
\cite{CCHM} --- cubes governed by the free De Morgan algebra, with
diagonals and reversal --- does not.  This paper is about what
replaces it.

Write $\cube$ for the De Morgan cube category: objects $[n]$, maps
$[m] \to [n]$ the $n$-tuples of elements of the free De Morgan
algebra $\DM_m$; equivalently, the opposite of the category of
finitely generated free De Morgan algebras.  Presheaves on $\cube$
are the De Morgan cubical sets interpreting CCHM type theory, and
the maps of $\cube$ are exactly the \emph{reparametrizations} that a
proof assistant's kernel applies silently: the \emph{strict layer}
of the theory, as opposed to the weak layer of Kan compositions.
A companion paper \cite{DMFibers} determined the boundary fibers of
$\DM_n$: the set of elements with prescribed faces is a Boolean
interval $2^{D}$, where $D$ is an antichain of \emph{diagonal}
points of the literal cube, with extremal fibers of size
$2^{\binom{n}{\lfloor n/2\rfloor}}$ by Sperner's theorem.  Thus a
boundary does not determine its filler --- degenerate-style cells
are non-unique in a quantified, Boolean-graded way.  The present
paper turns that local computation into a global structure theory
for $\cube$, in three steps.

\subparagraph*{The Karoubi envelope
  (\S\ref{sec:karoubi}).}  First, the negative fact in its sharpest
form: $\cube$ is not idempotent complete --- the idempotent
$x \mapsto x \vee \neg x$ on $[1]$ does not split
(Proposition~\ref{prop:nonsplit}), so $\cube$ is not a generalized
Reedy category and in particular not Eilenberg--Zilber, by the same
first obstruction Campion \cite{Campion} exhibited for the
diagonal--connection sites.  Campion's classification, however,
never meets $\cube$: his cube categories are categories of
$\{0,1\}$-valued operations, so his full-signature site is the
\emph{Boolean} cube category, where $x \vee \neg x = 1$.  We
identify the Karoubi envelope of the genuinely De Morgan site:
\[
  \Kar(\cube) \;\simeq\;
  (\text{finite projective De Morgan algebras})^{\mathrm{op}},
\]
and finite projective De Morgan algebras were characterized
combinatorially by Bova and Cabrer \cite{BovaCabrer} in the theory
of unification: dually, involutive posets satisfying three explicit
conditions.  The two literatures do not cite each other; the
identification, though easy once stated, appears to be new, and
slots $\cube$ into the table of cubical Karoubi envelopes next to
Sattler's Dedekind computation \cite{SattlerNote} and Campion's
Boolean one.

\subparagraph*{The defect formula (\S\ref{sec:defect}).}  Second,
the positive replacement for degeneracy normal forms.  Every
$g \in \DM_n$ determines the antichain $A_g$ of its minimal
\emph{diagonal} points, and we prove
(Theorem~\ref{thm:defect}) that
\[
  \rmin(g) \;:=\; g \setminus A_g
\]
is, in a single step, the least element of the boundary fiber of
$g$.  Consequently, for any substitution $f$ and any
boundary-minimal $\varphi$, the failure of $f^{*}\varphi$ to remain
minimal is measured by the explicit antichain $A_{f^{*}\varphi}$ ---
the \emph{defect} --- and corrected by deleting it, a move that
stays inside one Boolean fiber and is therefore connected to the
uncorrected transport by the strict median cylinder of
\cite{DMFibers}.

\subparagraph*{The coherence dichotomy (\S\ref{sec:coherence}).}
Third, the main theorem.  Correcting commutes with transporting ---
$\rmin(f^{*}g) = \rmin(f^{*}\rmin(g))$ --- if and only if the
transported boundary never reads a defect point:

\begin{thmstar}[Theorem~\ref{thm:dichotomy}]
  For a substitution $f : [m] \to [n]$ with induced map $F$ of
  literal cubes, and $g \in \DM_n$, the following are equivalent:
  \begin{enumerate}
  \item $\rmin(f^{*}g) = \rmin(f^{*}\rmin(g))$;
  \item $f^{*}g$ and $f^{*}\rmin(g)$ have equal boundaries;
  \item $F(\Face(\lc_m)) \cap A_g = \emptyset$.
  \end{enumerate}
\end{thmstar}

Each implication has a two-line proof, and the trichotomy of
generators falls out: faces, symmetries, reversal (and the
square-degeneracies $x \mapsto x \wedge \neg x$) send face points to
face points, so on the wide subcategory they generate the minimal
skeleton is a strictly natural retraction; degeneracies, connections
and diagonals map some face points onto diagonals, and for them the
criterion is sharp.  Two consequences deserve emphasis.  The
exceptional set is genuinely nonempty already for
\emph{degeneracies}: the cylinder $g \times I$ has $g$ itself among
its faces, so the minimal skeleton is \textbf{not cylinder-stable}
--- a precise expression of how far the strict layer is from a
Reedy-style theory.  And the criterion is exact, not merely
sufficient: on all $7{,}828{,}354$ elements of $\DM_3$ the three
conditions of the theorem carve out the same $1{,}133{,}756$
failures for a diagonal and for a connection substitution
(\S\ref{sec:machine}).

\subparagraph*{Thinness as structure (\S\ref{sec:thin}).}  In
cubical $\omega$-categories with connections, thin fillers of
commutative shells exist \emph{uniquely} (Higgins \cite{Higgins},
Steiner \cite{Steiner}); the uniqueness depends on strict
composition and on the normal forms of an Eilenberg--Zilber
site, and Steiner anticipated that in weak settings thin fillers
would exist ``but not necessarily'' uniquely.  The De Morgan strict layer realizes that anticipation
in a free, composition-less setting: fillers form Boolean torsors,
no property singles one out functorially (the dichotomy), and the
canonical data that remains --- minimal filler, defect antichain,
median cylinder --- is \emph{structure}.  We take this as the
conceptual summary: \emph{where Eilenberg--Zilber fails, thinness
cannot be a property; the median-term algebra is the minimal
structure replacing it.}

\subparagraph*{Relation to other work.}  Buchholtz--Morehouse
\cite{BuchholtzMorehouse} proved that the De Morgan cube category is
a \emph{strict test category}, so the homotopy theory over $\cube$
is standard --- our subject is the strict layer that presents it,
which their analysis leaves untouched.  Cavallo--Sattler
\cite{CavalloSattler} eliminate reversals up to homotopy over
self-dual interval theories; their twist construction
$T(I) = I \times I$ is, on representables, exactly the literal-cube
presentation used here (the diagonal points being the twist
diagonal), and their results are complementary: elimination is
up-to-path and model-level, while the judgmental layer of
implemented proof assistants remains De Morgan.
Dor\'e--Cavallo--M\"ortberg \cite{DCM} search for single fillers
over the same algebras; the present results describe the whole
solution set and its (non-)functoriality.  All machine verification
was performed with the scripts accompanying \cite{DMFibers},
extended for this paper (\S\ref{sec:machine}).

\section{Preliminaries}
\label{sec:prelim}

\subsection{The De Morgan cube category}

$\DM_n$ denotes the free De Morgan algebra on $x_1, \dots, x_n$; we
freely use its classical presentation (see \cite{DMFibers} and the
references there): the lattice of monotone Boolean functions on the
\emph{literal cube} $\lc_n = \{0,1\}^{2n}$, coordinates grouped in
pairs $(x_i, \neg x_i)$, elements identified with up-sets.  A point
is a \emph{face point} if some pair equals $(0,1)$ or $(1,0)$, a
\emph{diagonal point} if every pair is $(0,0)$ or $(1,1)$; write
$\Face(\lc_n)$ and $\Diag(\lc_n)$.  The \emph{faces}
$\bnd_i^e : \DM_n \to \DM_{n-1}$ substitute $x_i := e$; the
\emph{boundary} $\bnd g$ is the tuple of all faces.

The \emph{De Morgan cube category} $\cube$ has objects $[n]$
($n \geq 0$) and morphisms $[m] \to [n]$ the $n$-tuples
$f = (t_1, \dots, t_n)$, $t_i \in \DM_m$, composing by
substitution; it is the opposite of the category of finitely
generated free De Morgan algebras, and contains faces,
degeneracies, symmetries, both connections
($x_i := x_i \wedge x_j$, $x_i := x_i \vee x_j$), reversal
($x_i := \neg x_i$), and diagonals ($x_i := x_j$).  A morphism $f$
induces the substitution $f^{*} : \DM_n \to \DM_m$ and the
\emph{literal-cube map} $F : \lc_m \to \lc_n$,
$F(\alpha)_i = \bigl(t_i(\alpha), (\neg t_i)(\alpha)\bigr)$, so that
$(f^{*}g)(\alpha) = g(F(\alpha))$.  Note that $F$ is a map of
\emph{all} points and need not preserve the face/diagonal
dichotomy; this freedom is where everything below happens.

\subsection{The fiber theorem}

We import from \cite{DMFibers}: every fiber of the boundary map on
$\DM_n$ is an interval $[m, M]$ whose difference
$D = M \setminus m$ is an antichain of diagonal points; the fiber
is Boolean, canonically $2^{D}$; every antichain of diagonal points
arises; each compatible sphere has a least filler (the up-closure of
its prescribed points); boundaries are in bijection with antichains
of face points; and for same-boundary $\varphi, \varphi'$ the median
interpolation
$\chi = (\varphi \wedge \neg k) \vee (k \wedge \varphi') \vee
(\varphi \wedge \varphi')$ is a cylinder in one fresh variable,
constant on all faces.  An element is \emph{boundary-minimal} if it
is the least element of its fiber, equivalently if all its minimal
points are face points.

\section{The Karoubi envelope}
\label{sec:karoubi}

\begin{prop}
  \label{prop:nonsplit}
  The endomorphism $e = (x \mapsto x \vee \neg x)$ of $[1]$ is
  idempotent and does not split.  Hence $\cube$ is not idempotent
  complete, not a generalized Reedy category, and not an
  Eilenberg--Zilber category.
\end{prop}

\begin{proof}
  Idempotency is the absorption
  $(x \vee \neg x) \vee \neg(x \vee \neg x)
  = (x \vee \neg x) \vee (\neg x \wedge x) = x \vee \neg x$.
  Suppose $e = s \circ p$ with $p : [1] \to [k]$, $s : [k] \to [1]$
  and $p \circ s = \mathrm{id}_{[k]}$.  Contravariantly
  $s^{*} \circ p^{*} = \mathrm{id}_{\DM_k}$, so $\DM_k$ is a retract
  of $\DM_1$; since $|\DM_1| = 6$ and $|\DM_k| = M(2k) \geq 168$ for
  $k \geq 2$, we must have $k \leq 1$.  For $k = 0$, $e$ factors
  through $[0]$, so $e^{*}$ is constant; but
  $x \vee \neg x \notin \{0, 1\}$ in $\DM_1$.  For $k = 1$, the
  retraction forces $s^{*}$ to be bijective (finite cardinality), so
  $e^{*} = p^{*} \circ s^{*}$ is injective; but
  $e^{*}(x) = x \vee \neg x = e^{*}(x \vee \neg x)$ while
  $x \neq x \vee \neg x$.  The last two claims follow since
  generalized Reedy categories are idempotent complete
  \cite{Campion}.
\end{proof}

\begin{thm}
  \label{thm:karoubi}
  The Karoubi envelope of the De Morgan cube category is the
  opposite of the category of finite projective De Morgan algebras:
  \[
    \Kar(\cube) \;\simeq\;
    (\mathsf{DM}\text{-}\mathsf{proj}_{\mathrm{fin}})^{\mathrm{op}}.
  \]
  Its objects are characterized combinatorially by Bova--Cabrer
  \cite[Thm.~11]{BovaCabrer}: under Cornish--Fowler duality
  \cite{CornishFowler}, a finite De Morgan algebra is projective iff
  its dual involutive poset $(P, \leq, i)$ is a nonempty lattice, every
  $x \leq i(x)$ admits $y \geq x$ with $y = i(y)$, and the sub-poset
  $\{x \mid x \leq i(x)\}$ is $3$-complete.
\end{thm}

\begin{proof}
  $\cube$ is the opposite of the finitely generated free De Morgan
  algebras.  In any variety, idempotents split in the category of
  algebras, and retracts of free algebras are exactly the
  projectives; hence the idempotent completion of the full
  subcategory of f.g.\ frees is the full subcategory of
  f.g.\ projectives.  The variety of De Morgan algebras is locally
  finite, so finitely generated projectives are finite
  \cite{BovaCabrer}.  Taking opposites gives the displayed
  equivalence, and the characterization of the objects is
  Bova--Cabrer's theorem.
\end{proof}

\begin{rem}
  This places $\cube$ in the table of cubical Karoubi envelopes:
  Dedekind cubes with diagonals complete to finite lattices with
  monotone maps (Sattler \cite{SattlerNote}, cf.\
  \cite[Prop.~8.8]{Campion}); Boolean cubes to finite nonempty sets
  \cite[Prop.~8.11]{Campion}; De Morgan cubes to the dual of finite
  projective De Morgan algebras.  The unification-theoretic
  literature where the projectivity characterization lives
  \cite{BovaCabrer} and the cubical literature appear never to have
  cited one another.  The same argument identifies the Karoubi
  envelope of the \emph{Kleene} cube category via
  \cite[Thm.~12]{BovaCabrer}.  We note that Campion's negative
  results \cite[Thm.~8.12(2)]{Campion} make it unlikely that any
  degree function turns these envelopes into Eilenberg--Zilber
  categories; the sequel does not attempt this, and instead
  isolates what survives.
\end{rem}

\section{Minimal fillers: the defect formula}
\label{sec:defect}

For $g \in \DM_n$ let $A_g$ be the set of minimal points of $g$
(qua up-set) that are diagonal.  Minimal points form an antichain,
so $A_g$ is an antichain of diagonal points.

\begin{thm}[defect formula]
  \label{thm:defect}
  For every $g \in \DM_n$, the element
  $\rmin(g) := g \setminus A_g$ is the least element of the
  boundary fiber of $g$.  In particular one deletion step reaches
  the fiber minimum, and $g$ is boundary-minimal iff
  $A_g = \emptyset$.
\end{thm}

\begin{proof}
  $\rmin(g)$ is an up-set: a point strictly above $\alpha \in A_g$
  is either a face point (hence in $g$, hence in $\rmin(g)$), or a
  diagonal point $\beta > \alpha$, in which case the companion
  lemma \cite[Lem.~2.3]{DMFibers} provides a face point $\gamma$
  with $\alpha < \gamma \leq \beta$; then $\gamma \in g$, so
  $\beta \in \Up\gamma \subseteq \rmin(g)$.  Its boundary equals
  $\bnd g$ since only diagonal points were removed.  No new
  diagonal minimal points appear: if $\beta$ is diagonal and minimal
  in $\rmin(g)$, then $\beta$ was not minimal in $g$ (else
  $\beta \in A_g$), so some $\delta < \beta$ lies in $g$; taking
  $\delta$ minimal in $g$, either $\delta$ is a face point --- but
  then $\delta \in \rmin(g)$, contradicting minimality of $\beta$
  --- or $\delta \in A_g$ is diagonal, and the companion lemma gives
  a face point $\gamma \in (\delta, \beta]$, with
  $\gamma \neq \beta$ ($\beta$ is diagonal), lying in $\rmin(g)$ ---
  again a contradiction.  So all minimal points of $\rmin(g)$ are
  face points, which characterizes the least element of the fiber
  \cite[Thm.~4.2]{DMFibers}.
\end{proof}

\begin{cor}[the defect of a substitution]
  \label{cor:defect}
  Let $f : [m] \to [n]$ be a morphism of $\cube$ and
  $\varphi \in \DM_n$ boundary-minimal.  Then the deviation of
  $f^{*}\varphi$ from boundary-minimality is the explicit antichain
  \begin{multline*}
    A_{f^{*}\varphi} \;=\;
    \{\, \alpha \in \Diag(\lc_m) \;:\;
    F(\alpha) \in \varphi \text{ and}\\
    F(\beta) \notin \varphi
    \text{ for every lower cover } \beta \lessdot \alpha \,\},
  \end{multline*}
  and $\rmin(f^{*}\varphi)$ differs from $f^{*}\varphi$ exactly by
  this antichain, inside one Boolean fiber; the median cylinder of
  \cite{DMFibers} connects the two by a strict homotopy constant on
  the boundary.
\end{cor}

\section{The coherence dichotomy}
\label{sec:coherence}

The retraction $\rmin$ is a normal form for \emph{fixed} $[n]$; the
question is its interaction with substitution.  Call
$f : [m] \to [n]$ and $g \in \DM_n$ \emph{coherent} if
$\rmin(f^{*}g) = \rmin(f^{*}\rmin(g))$.

\begin{thm}[coherence dichotomy]
  \label{thm:dichotomy}
  For every morphism $f$ of $\cube$ and every $g \in \DM_n$, the
  following are equivalent:
  \begin{enumerate}
  \item $(f, g)$ is coherent;
  \item $\bnd(f^{*}g) = \bnd(f^{*}\rmin(g))$;
  \item $F(\Face(\lc_m)) \cap A_g = \emptyset$.
  \end{enumerate}
\end{thm}

\begin{proof}
  (3)$\Rightarrow$(2): $g$ and $\rmin(g)$ agree outside $A_g$, and
  the boundaries of the transports only read the values of $g$,
  resp.\ $\rmin(g)$, on $F(\Face(\lc_m))$.
  (2)$\Rightarrow$(1): elements with equal boundaries lie in one
  fiber, whose least element is unique.
  (1)$\Rightarrow$(3), contrapositively: if $\alpha$ is a face point
  with $F(\alpha) \in A_g$, then
  $(f^{*}g)(\alpha) = g(F(\alpha)) = 1$ while
  $(f^{*}\rmin g)(\alpha) = 0$; a disagreement at a face point is a
  disagreement of boundaries, and least elements determine their
  boundaries ($\rmin(a) = \rmin(b)$ forces
  $\bnd a = \bnd \rmin(a) = \bnd \rmin(b) = \bnd b$), so the two
  corrected elements differ.
\end{proof}

\begin{prop}[classification of generators]
  \label{prop:generators}
  Faces, symmetries, reversal, and the substitutions
  $x_i := x_i \wedge \neg x_i$, $x_i := x_i \vee \neg x_i$ send face
  points to face points; hence they are coherent for \emph{every}
  $g$, and the wide subcategory $\cubefp \subseteq \cube$ of
  face-point-preserving morphisms (which they generate, and which is
  closed under composition) carries $\rmin$ as a strictly natural
  retraction onto boundary-minimal elements.  Degeneracies,
  connections, and diagonals each map some face point to a diagonal
  point, and for each of them elements $g$ violating
  condition~(3) exist.
\end{prop}

\begin{proof}
  For a face $[n-1] \to [n]$, $F$ embeds $\lc_{n-1}$ into a face
  subcube, sending every point (facial or not) to a face point.
  Symmetries and reversal permute pair-coordinates and swap pair
  entries, preserving the dichotomy.  For
  $x_i := x_i \wedge \neg x_i$: if $\alpha$ is facial via pair $j
  \neq i$, so is $F(\alpha)$; if via pair $i$, then
  $(x_i \wedge \neg x_i, x_i \vee \neg x_i)$ evaluates at
  $(0,1)$ or $(1,0)$ to $(0, 1)$, a facial pair.  Closure under
  composition is immediate ($F' F(\Face) \subseteq F'(\Face)
  \subseteq \Face$), and naturality on $\cubefp$ is
  Theorem~\ref{thm:dichotomy}(3$\Rightarrow$1) applied to all $g$.
  For a degeneracy $[n+1] \to [n]$, a point facial only via the
  dropped pair maps to a diagonal point; for the connection
  $x_1 := x_1 \wedge x_2$, the point with pairs
  $((1,0), (0,0), \dots)$ maps to $((0,0), (0,0), \dots)$; for the
  diagonal $x_1 := x_2$, the point $((0,1), (0,0), \dots)$ maps to
  $((0,0), (0,0), \dots)$.  Violating elements exist in each case:
  take $g$ with $A_g$ containing the reached diagonal point
  (realizability in \cite[Thm.~3.4]{DMFibers}).
\end{proof}

\begin{rem}[the skeleton is not cylinder-stable]
  \label{rem:cylinder}
  The degeneracy case deserves emphasis.  For the cylinder
  $g \times I := \sigma^{*} g$ along $\sigma : [n+1] \to [n]$, the
  two new faces are $g$ itself; if $A_g \neq \emptyset$ then
  condition~(3) fails (every diagonal point of $\lc_n$ is the image
  of a point facial only via the dropped pair), and indeed
  $\rmin(g \times I) \neq \rmin(g) \times I$: the minimal filler of
  the cylinder's boundary is not the cylinder of the minimal
  filler.  A Reedy-style ``skeleton'' functor is thus unavailable in
  principle --- consistent with, and sharpening,
  Proposition~\ref{prop:nonsplit}.
\end{rem}

\section{Thinness as structure}
\label{sec:thin}

In cubical $\omega$-categories with connections, Higgins
\cite{Higgins} proved that every commutative shell has a
\emph{unique} thin filler, and Steiner \cite{Steiner} that the same
uniqueness organizes the cubical nerves of $\omega$-categories;
in both cases ``thin'' is definable from the ambient composition
structure (Higgins: folding into an identity; and thin structures
are equivalent to connections, \cite[Thm.~3.1]{Higgins}).  Steiner
remarked that for weak higher categories one expects thin fillers
``to exist but not necessarily to be unique''.  The De Morgan
strict layer is an extreme test case: there is no composition at
all, fillers of a boundary form the Boolean torsor $2^{D}$
\cite{DMFibers}, and the dichotomy
(Theorem~\ref{thm:dichotomy}) shows that no choice of
representatives is stable under all reparametrizations --- the
obstruction being reached by degeneracies, connections and
diagonals alike.  What does survive is \emph{structure}: the
minimal filler ($\rmin$), the defect antichain
(Corollary~\ref{cor:defect}) measuring every transport's deviation,
and the median cylinder correcting it within the fiber.  In slogan
form: \emph{without Eilenberg--Zilber normal forms, thinness cannot
be a property; the minimal-filler-plus-median package is the minimal
structure that replaces it.}  We expect this package, rather than
any uniqueness principle, to be the right input for algebraic
weak-factorization-system treatments of the strict layer; we leave
that organization to future work
(\S\ref{sec:outlook}).

\section{Machine verification}
\label{sec:machine}

All theorems above are elementary, but their statements were found,
and are guarded, by exhaustive machine checks (scripts accompany the
artifact; see \cite{DMFibers} for the census infrastructure).
Through $n \leq 3$:
the defect formula of Theorem~\ref{thm:defect} was verified on all
$168$ elements of $\DM_2$ and all $7{,}828{,}354$ elements of
$\DM_3$ against brute-force fiber minima (zero discrepancies); the
coherence dichotomy was verified on all of $\DM_2$ for five
generating substitutions ($81$ incoherent instances, matching
$23 + 23 + 24 + 11$ counted independently), and on all of $\DM_3$
for a diagonal, a connection, and the reversal --- conditions
(1), (2), (3) of Theorem~\ref{thm:dichotomy} carve out identical
subsets, of sizes $1{,}133{,}756$, $1{,}133{,}756$, and $0$
respectively.  At $n = 2$ the exceptional set for connections and
diagonals consists exactly of the two elements whose defect is a
pole of the literal cube ($g = 1$ and the four-fold conjunction
$x \wedge \neg x \wedge y \wedge \neg y$); from $n = 3$ on it is
large --- about $14.5\%$ of all elements for the diagonal
substitution --- so the dichotomy is a genuinely quantitative
constraint, not an edge-case bookkeeping device.

\section{Outlook}
\label{sec:outlook}

Three directions.  (1) \emph{Algebraic organization}: package
(minimal filler, defect, median correction) as an algebraic weak
factorization system on De Morgan cubical sets in which the
boundary inclusions' right lifting structures are the canonical
minimal fillers and the left class carries the Boolean fiber data;
the dichotomy says any such structure will be genuinely
\emph{algebraic} over $\cube \setminus \cubefp$, never property-like
--- the verification of the distributivity axioms against the
$\chi$-corrections is the remaining work.  (2) \emph{The Karoubi
envelope as a site}: Theorem~\ref{thm:karoubi} invites a study of
presheaves on finite projective De Morgan algebras --- the
``involutive lattices'' picture of \cite{BovaCabrer} suggests
comparisons both with Sattler's finite-lattice site and with the
twist construction of \cite{CavalloSattler}, under which our
literal cube is the twist of the Dedekind evaluation cube.
(3) \emph{Kleene and beyond}: every argument here uses only the
literal-cube combinatorics; the Kleene case constrains but does not
excise the diagonal points, and the corresponding fiber, defect and
coherence theory (with \cite[Thm.~12]{BovaCabrer} supplying the
Karoubi envelope) should interpolate between the De Morgan and
Dedekind extremes, where fibers are Boolean of positive rank,
respectively trivial.

\subparagraph*{Artifact.}  The verification scripts and the
machine-checked type-theoretic counterparts of the companion
results are available in the repository at
\url{https://github.com/r-shibusawa/cubical-strict-j}
(archived at DOI \href{https://doi.org/10.5281/zenodo.21405961}%
{10.5281/zenodo.21405961}).

\begin{thebibliography}{10}

\bibitem{BovaCabrer}
S.~Bova, L.~Cabrer.
\newblock Unification and projectivity in De Morgan and Kleene
  algebras.
\newblock \emph{Order}\ 31 (2014), 159--187.

\bibitem{BuchholtzMorehouse}
U.~Buchholtz, E.~Morehouse.
\newblock Varieties of cubical sets.
\newblock \emph{RAMiCS 2017}, LNCS 10226, 2017.

\bibitem{Campion}
T.~Campion.
\newblock Cubical sites as Eilenberg--Zilber categories.
\newblock Preprint, arXiv:2303.06206, 2023.

\bibitem{CavalloSattler}
E.~Cavallo, C.~Sattler.
\newblock Eliminating reversals from cubical type theories.
\newblock \emph{LICS 2026}, LIPIcs 380, 2026.

\bibitem{CCHM}
C.~Cohen, T.~Coquand, S.~Huber, A.~M\"ortberg.
\newblock Cubical type theory: a constructive interpretation of the
  univalence axiom.
\newblock \emph{TYPES 2015}, LIPIcs 69, 2018.

\bibitem{CornishFowler}
W.~H.~Cornish, P.~R.~Fowler.
\newblock Coproducts of De Morgan algebras.
\newblock \emph{Bull.\ Austral.\ Math.\ Soc.}\ 16 (1977), 1--13.

\bibitem{DCM}
M.~Dor\'e, E.~Cavallo, A.~M\"ortberg.
\newblock Automating boundary filling in cubical type theories.
\newblock \emph{Log.\ Methods Comput.\ Sci.}\ 22(2):28, 2026.

\bibitem{Higgins}
P.~J.~Higgins.
\newblock Thin elements and commutative shells in cubical
  $\omega$-categories.
\newblock \emph{Theory Appl.\ Categ.}\ 14 (2005), 60--74.

\bibitem{SattlerNote}
C.~Sattler.
\newblock Idempotent completion of cubes in posets.
\newblock Preprint, arXiv:1805.04126, 2019.

\bibitem{DMFibers}
R.~Shibusawa.
\newblock Boundary fibers in free De Morgan algebras: Boolean
  intervals, Sperner bounds, and median contractions.
\newblock Preprint, 2026.  Included in the artifact repository
  \url{https://github.com/r-shibusawa/cubical-strict-j} (from
  release \texttt{v1.6.0}, DOI
  \href{https://doi.org/10.5281/zenodo.21695307}%
  {10.5281/zenodo.21695307}).

\bibitem{Steiner}
R.~Steiner.
\newblock Thin fillers in the cubical nerves of omega-categories.
\newblock \emph{Theory Appl.\ Categ.}\ 16 (2006), 144--173.

\end{thebibliography}

\end{document}
